\section{The next construction problem}
\label{qdn-next-section}

The normalized local quantum map and its finite-history compatibility
now have a specified construction. Their domains and positive
majorants are uniform in preceding homogeneous classical depth.
They provide the setting in which an equivalent connected
representation can be used at the next local integration.
Field-independent vacuum factors remain separately tracked;
no convergence of an extensive vacuum energy is inferred.

The actual inverse coupling obeys
\[
 \gamma_{j+1}=\gamma_j-
       \beta_{j+1}\bigl(\mathcal F_{j,g_j}(E_j)\bigr),
 \qquad \gamma_j=g_j^{-2}.
\]
Completing its weak-coupling trajectory is the next essential
step. Absolute increment bounds alone do not establish its
direction or permit the ultraviolet cutoff to be removed at a
fixed interacting endpoint. The generated drift must be evaluated
with the actual determinants, cutoff factors and inherited
interaction. Its normalization must agree with the reference
action used in the exact density identity.

The activity estimates concern the specified homogeneous
small-field sector. The full Wilson construction must also
incorporate mixed and large-field histories, identify the
continued local representation on their required gauge-orbit
domains, and obtain uniform estimates with observable sources.
Those estimates feed the compactness, rotation, nontriviality
and reconstruction obligations in Notebook 01.

\begin{thebibliography}{9}
\bibitem{priorqdn}
\emph{Toward a four-dimensional Yang--Mills construction},
research Notebooks 15--29 in this workspace. The ordinary minimum
and gauge transfer, strict overlapping Gaussian expansion,
common-domain quantum map and total Ward extraction are the
principal inputs.
\bibitem{balaban109qdn}
T.~Ba\l aban,
\emph{Renormalization group approach to lattice gauge field
theories. I. Generation of effective actions in a small field
approximation and a coupling constant renormalization in four
dimensions},
Commun. Math. Phys. \textbf{109} (1987), 249--301.
\href{https://doi.org/10.1007/BF01215223}{Publisher record}.
\bibitem{balaban116qdn}
T.~Ba\l aban,
\emph{Renormalization group approach to lattice gauge field
theories. II. Cluster expansions},
Commun. Math. Phys. \textbf{116} (1988), 1--22.
\href{https://doi.org/10.1007/BF01239022}{Publisher record}.
The finite localization and resummation architecture is checked
against the original article. The quotient and local-continuity
proofs in this notebook are stated and proved separately.
\bibitem{bkarqdn}
A.~Abdesselam, notes on the BKAR formula, Theorem 1.
\href{https://mabdesselam.github.io/BKAR.pdf}{Author-hosted notes}.
The finite forest identity is used with the explicit positive
majorants and strict margins of the preceding construction.
\bibitem{sourceqdn}
E.~Sussman and collaborators,
\href{https://github.com/ethanwsussman/Yang--Mills-LLM-experiment}
{Yang--Mills-LLM-experiment}, deposited raw accounts at commit
\texttt{e392ef45013a581664bab76b90867170b30fe10e}.
\end{thebibliography}
\end{document}
